\section{Resultados esperados}

Ao final deste estudo, espera-se alcançar um modelo preditivo robusto que seja capaz de prever com precisão o comportamento de compra dos clientes em um ambiente de e-commerce, especificamente o tempo até a próxima compra ou o risco de abandono. O produto final será uma análise detalhada dos fatores que influenciam essas decisões de compra, com base em variáveis como frequência de visitas, histórico de compras e características demográficas dos clientes. Com isso, será possível identificar segmentos de clientes com maior probabilidade de realizar compras e aqueles em risco de abandonar a plataforma, permitindo que estratégias de marketing sejam mais eficazmente direcionadas.

Além disso, espera-se que os resultados possam fornecer insights valiosos para empresas de e-commerce, permitindo a otimização de campanhas de marketing, a personalização da experiência do usuário e a redução da rotatividade de clientes. O modelo de análise de sobrevivência aplicado, com ênfase no modelo de Cox, será uma ferramenta útil para que as empresas possam antecipar comportamentos e tomar decisões informadas para aumentar a fidelização e maximizar a rentabilidade.

Porém, algumas limitações precisam ser consideradas. Em primeiro lugar, os dados utilizados são sintéticos e podem não refletir perfeitamente a complexidade e a variabilidade dos comportamentos reais dos clientes em plataformas de e-commerce. Isso pode limitar a generalização dos resultados para outros contextos ou mercados. Além disso, o tempo disponível para a realização da pesquisa pode restringir a profundidade da análise, especialmente em relação à exploração de diferentes técnicas de modelagem ou à inclusão de outras variáveis que poderiam enriquecer o modelo preditivo.

Apesar dessas limitações, espera-se que os resultados obtidos contribuam para o aprimoramento do conhecimento científico na área de previsão de comportamento de clientes em e-commerce, especialmente no que tange à utilização de modelos de análise de sobrevivência em contextos digitais. Os insights gerados também têm o potencial de servir como uma fonte útil de informação para empresas de e-commerce, que poderão aplicar as recomendações para melhorar suas estratégias de marketing e aumentar a retenção de clientes, resultando em uma gestão mais eficiente e assertiva de seus recursos.
